\documentclass[12pt]{article}
\usepackage[utf8]{inputenc}
\usepackage{amsmath, amssymb}
\usepackage{booktabs}
\usepackage{geometry}
\usepackage{graphicx}
\usepackage{caption}
\geometry{margin=1in}

\title{Precertamen 3}
\author{}
\date{}

\begin{document}

\maketitle

\section*{Ejercicio 1}
Un investigador estima el siguiente modelo para explicar el ingreso anual ($Salario$) de individuos en una muestra aleatoria de 500 trabajadores residentes en una ciudad latinoamericana:

\begin{equation}
Salario_i = \beta_0 + \beta_1 \cdot Educacion_i + \beta_2 \cdot Experiencia_i + \beta_3 \cdot Genero_i + \beta_4 \cdot SectorPublico_i + \varepsilon_i
\end{equation}

\textbf{Variables}:
\begin{itemize}
  \item $Salario$: ingreso anual en dólares
  \item $Educacion$: años de escolaridad
  \item $Experiencia$: años de experiencia laboral
  \item $Genero$: dummy = 1 si es hombre, 0 si es mujer
  \item $SectorPublico$: dummy = 1 si trabaja en el sector público
  \item $Habilidad$: resultado en test estandarizado de habilidad laboral
  \item $DistanciaEscuela$: distancia (km) al colegio más cercano en la infancia
  \item $Hermanos$: número de hermanos
\end{itemize}

\section*{Estimación por MCO}

\begin{center}
\begin{tabular}{lccc}
\toprule
Variable & Coeficiente & Error Estándar & Valor-p \\
\midrule
Intercepto       & 10450*** & 1720 & 0.000 \\
Educacion        & 820***   & 110  & 0.000 \\
Experiencia      & 520***   & 85   & 0.000 \\
Genero           & 3100***  & 950  & 0.001 \\
SectorPublico    & 2100***  & 740  & 0.004 \\
\midrule
$R^2$ = 0.42 & & & \\
\bottomrule
\end{tabular}
\end{center}

\section*{Preguntas}

\subsection*{1. Variables Dummy}
\begin{enumerate}
  \item[a)] ¿Cómo se interpreta el coeficiente de $Genero$? ¿Qué sugiere este resultado?
  \item[b)] ¿Qué precauciones metodológicas deben considerarse al incluir variables dummy?
\end{enumerate}

\subsection*{2. Multicolinealidad}

\textbf{Factores de inflación de varianza (VIF)}:

\begin{center}
\begin{tabular}{lc}
\toprule
Variable & VIF \\
\midrule
Educacion       & 6.4 \\
Experiencia     & 5.8 \\
Genero          & 1.4 \\
SectorPublico   & 1.5 \\
\bottomrule
\end{tabular}
\end{center}

\begin{enumerate}
  \item[a)] ¿Existe multicolinealidad? ¿Qué implicancias tiene para los estimadores de MCO?
  \item[b)] ¿Qué estrategias puede seguir el investigador para mitigar este problema?
\end{enumerate}

\subsection*{3. Endogeneidad y Variables Instrumentales}

Se sospecha que $Educacion$ está correlacionada con el término de error. Se proponen dos instrumentos:
\begin{itemize}
  \item $DistanciaEscuela$: distancia al colegio más cercano en la infancia
  \item $Hermanos$: número de hermanos
\end{itemize}

\textbf{Primera etapa (regresión auxiliar)}:
\begin{equation*}
Educacion_i = \pi_0 + \pi_1 \cdot DistanciaEscuela_i + \pi_2 \cdot Hermanos_i + \pi_3 \cdot Experiencia_i + \pi_4 \cdot Genero_i + \pi_5 \cdot SectorPublico_i + v_i
\end{equation*}

\begin{center}
\begin{tabular}{lcc}
\toprule
Variable & Coeficiente & Error Estándar \\
\midrule
DistanciaEscuela & -1.20 & (0.40) \\
Hermanos         & -0.35 & (0.15) \\
\midrule
$R^2$ = 0.31 & & \\
\bottomrule
\end{tabular}
\end{center}

\textbf{Segunda etapa (2SLS)}:
\begin{equation*}
Salario_i = \beta_0 + \hat{\beta}_1 \cdot EducacionInstrumentada_i + \beta_2 \cdot Experiencia_i + \beta_3 \cdot Genero_i + \beta_4 \cdot SectorPublico_i + \varepsilon_i
\end{equation*}

\begin{center}
\begin{tabular}{lcc}
\toprule
Variable & Coeficiente & Error Estándar \\
\midrule
Educacion        & 1150  & (250) \\
Experiencia      & 400   & (95)  \\
Genero           & 2900  & (980) \\
SectorPublico    & 2000  & (780) \\
\midrule
$R^2$ = 0.39 & & \\
\bottomrule
\end{tabular}
\end{center}

\begin{enumerate}
  \item[a)] ¿Por qué $Educacion$ podría ser endógena en este modelo?
  \item[b)] Compare la validez de los instrumentos propuestos. ¿Cuál cumple mejor con los requisitos?
  \item[c)] Compare los coeficientes de $Educacion$ entre MCO y 2SLS. ¿Qué interpretación tiene esta diferencia?
\end{enumerate}

\subsection*{4. Heterocedasticidad}

Se realiza el test de Breusch-Pagan y se obtiene un p-valor = 0.004.
\begin{enumerate}
  \item[a)] ¿Qué indica este resultado sobre los supuestos del modelo?
  \item[b)] ¿Qué consecuencias tiene sobre los errores estándar y la inferencia estadística?
  \item[c)] ¿Cómo puede corregirse este problema?
\end{enumerate}

\subsection*{5. Reflexión Final}
Analiza cómo cada uno de los problemas discutidos (dummies, multicolinealidad, endogeneidad y heterocedasticidad) afecta la validez del modelo. ¿Cuál de ellos consideras más grave en este contexto y por qué?

\end{document}
